\section{Aneks replikacyjny}

Eksperyment wykonano w lokalnym worktree repozytorium \amor{} na gałęzi \path{mba/basic-income-scenario}. Punktem odniesienia był produkcyjny baseline modelu, skalibrowany do snapshotu gospodarki Polski na dzień 2026-04-30. Zakres i status tej powierzchni baseline są udokumentowane w następujących artefaktach repozytorium:

\begin{itemize}
    \item \path{docs/calibration-register.md} -- rejestr parametrów, jednostek, źródeł, transformacji i statusów kalibracji; wskazuje, że produkcyjne ustawienia baseline odnoszą się do snapshotu Polski z 2026-04-30.
    \item \path{docs/scenario-registry.md} -- rejestr scenariuszy; definiuje scenariusz baseline jako punkt odniesienia dla kontrfaktów oraz wskazuje jego źródło i datę jako produkcyjny baseline Polski 2026-04-30.
    \item \path{docs/empirical-validation-report.md} oraz \path{docs/empirical-validation/} -- raport i artefakty walidacyjne pokazujące, jak outputy modelu są porównywane z zewnętrznymi targetami empirycznymi.
    \item \path{docs/mba-basic-income-sweep-plan.md} -- plan eksperymentu BDP; dokumentuje, że BDP jest scenariuszem polityki publicznej względem baseline, a nie elementem kalibracji stanu początkowego.
\end{itemize}

Scenariusz BDP uruchamiano przez standardowy plik JAR po wykonaniu \texttt{sbt assembly}. Każdy poziom BDP był odrębnym uruchomieniem w wariancie centralnym \(\lambda = 0{,}5\):

\begin{verbatim}
sbt assembly
java -jar target/scala-3.8.2/amor-fati.jar 10 robust-bdp-0000 \
  --duration 60 --run-id robust-bdp-0000-60m-10s --bdp 0
java -jar target/scala-3.8.2/amor-fati.jar 10 robust-l050-bdp-0500 \
  --duration 60 --run-id robust-l050-bdp-0500-60m-10s \
  --bdp 500 --bdp-lambda 0.5
...
java -jar target/scala-3.8.2/amor-fati.jar 10 robust-l050-bdp-3000 \
  --duration 60 --run-id robust-l050-bdp-3000-60m-10s \
  --bdp 3000 --bdp-lambda 0.5
\end{verbatim}

Pliki wynikowe zapisano w katalogu \path{mc/}. Dla każdego poziomu powstało dziesięć plików seed CSV oraz agregaty \path{hh}, \path{banks} i \path{firms}. Przykładowy wzorzec nazwy pliku:

\begin{verbatim}
mc/robust-l050-bdp-2000_robust-l050-bdp-2000-60m-10s_60m_seed001.csv
\end{verbatim}

Analizę odporności kanału płacy rezerwowej uruchomiono dla \(\lambda = 0\), \(0{,}25\), \(0{,}5\), \(0{,}75\) i \(1\), przy tym samym horyzoncie 60 miesięcy. Każda kombinacja \(\lambda\) i dodatniego poziomu BDP była osobnym uruchomieniem JAR-a, po 10 seedów:

\begin{verbatim}
java -jar target/scala-3.8.2/amor-fati.jar 10 robust-l075-bdp-2000 \
  --duration 60 --run-id robust-l075-bdp-2000-60m-10s \
  --bdp 2000 --bdp-lambda 0.75
\end{verbatim}

Agregację wyników, tabele terminalne i rysunki generują skrypty:
\begin{verbatim}
paper/analysis/figures_bdp_sweep.py
paper/analysis/figures_lambda_robustness.py
\end{verbatim}

Skrypty zapisują tabele \path{paper/latex/data/bdp_sweep_terminal_summary.csv} i \path{paper/latex/data/bdp_lambda_robustness_terminal_summary.csv} oraz rysunki w katalogu \path{paper/latex/figures/}.

Wersja finalna pracy powinna uzupełnić niniejszy aneks o hash commita i dokładną wersję JAR.
